\documentclass[10pt,a4paper]{amsart}
\usepackage[margin=19mm]{geometry}
\usepackage{amsmath,amssymb,mathtools}
\usepackage[scr=rsfs,cal=euler]{mathalfa}
\usepackage{xfrac}
\usepackage[unicode,hidelinks,bookmarks=false]{hyperref}
\newcommand{\ie}{i.e.\ }
\newcommand{\cf}{cf.\ }
\newcommand{\resp}{respectively\ }
\newcommand{\Subsection}{\subsection}
\providecommand{\nobreakdash}{\nobreak\hbox{-}}
\setlength{\emergencystretch}{2em}
\multlinegap0pt
\usepackage{amssymb}
\usepackage{xcolor}
\usepackage{float}
\newtheorem{theorem}{Theorem}
\newtheorem{definition}{Definition}
\title{On some Sobolev and P\'olya-Szeg\"o type inequalities with weights and applications}
\author{Trung Hieu Giang, Nguyen Minh Tri, Dang Anh Tuan}
\date{Source: arXiv:2412.15490v4 (CC BY 4.0)}
\begin{document}
\maketitle

\noindent\emph{Source and licence.} On some Sobolev and P\'olya-Szeg\"o type inequalities with weights and applications. Version arXiv:2412.15490v4 (CC BY 4.0), distributed by arXiv under the Creative Commons Attribution 4.0 International licence, http://creativecommons.org/licenses/by/4.0/, as recorded in the arXiv licence metadata for this exact version and shown on the abstract page https://arxiv.org/abs/2412.15490v4. Full licence notice: \texttt{licenses/arxiv-2412.15490-CC-BY-4.0.txt}.\par\smallskip

\noindent\textit{Editorial scope.} Counted unit: the article's theorem thm1 (source0.tex lines 106--112). The article's complete original proof of it is delivered with it (source0.tex lines 213--252, one \texttt{proof} environment of 40 lines, which the article places away from the statement). No mathematics is written, repaired or re-derived: the statement and the proof are the article's own lines, in the article's own notation, and no retained line is substituted or re-flowed.  Retained uncounted as the source-local context the proof needs: lines 186--189; lines 197--200.  Nothing was omitted from any retained span, so the package carries no omission record.  The retained lines cite 1 works, whose entries are transcribed verbatim from the article's own bibliography source in the e-print and delivered with short editorial labels recorded in the item's reference mapping.  No retained line contains a figure environment, an image-inclusion command or drawing code, so no image is delivered and the package declares no asset dependency.  Editorial formatting only, in addition: an AMS \texttt{amsart} preamble that restates the article's own declarations and macros used by the retained lines, namely the environments \texttt{theorem}, \texttt{definition} on separate counters, together with the standard packages those lines need.  The retained environments print in this document as Theorem 1, Definition 1 in order of appearance, whereas the article prints the same items with the numbering of its own full document.  The complete native source of the article as distributed in the arXiv e-print for this exact version is bundled unchanged as \texttt{sources/170444/source0.tex}.
	\begin{definition}\label{def1}
		The $(2,\alpha)$-volume of $E$ is defined by
		$$ |E|_{2,\alpha} = \int\limits_E |x|^{2\alpha} dx_1dx_2 dy.$$ 
	\end{definition}

	\begin{definition}\label{def3}
		For each $j=1,...,2n(\alpha)$, the $(2,\alpha,j)$-area of $E$ is defined by
		$$ P_{2,\alpha,j} (E) = \int\limits_{\partial_{s_j} E} |x|^\alpha \sqrt{ \nu_1^2 + \nu_2^2 + |x|^{2\alpha} \nu_3^2}d\mathcal{H}^2.$$ 
	\end{definition}

	\begin{theorem}[The weighted isoperimetric inequality]\label{thm1}
		Let $E \subset \mathbb{R}^3$ be a nonempty bounded open set with Lipschitz boundary $\partial E$ such that $P_{2,\alpha} (E)$ is finite. Then the following inequality holds for every $j=1,...,2n(\alpha)$
		\begin{equation}\label{eqn2}
			\dfrac{P_{2,\alpha,j}(B_{s_j})^{3/2}}{|B_{s_j}|_{2,\alpha}} \leq \dfrac{P_{2,\alpha}(E)^{3/2}}{|E|_{2,\alpha}},
		\end{equation}
		where $B_{s_j} := \left\{ (x_1,x_2,y) \in \mathbb{R}^3_{s_j}, \textrm{ }\dfrac{|x|^{2\alpha+2}}{(\alpha+1)^2} + y^2 < 1 \right\}$. Here, the notations $P_{2,\alpha}$, $P_{2,\alpha,j}$, and $|\cdot|_{2,\alpha}$ respectively denote weighted areas and weighted volume, which will be precisely given in Definitions \ref{def1}-\ref{def3}. 
	\end{theorem}

	\begin{proof}[Proof of Theorem \ref{thm1}]
		For readers' convenience, we will divide the proof into two steps.
		
		\noindent\textbf{Step 1:} Assuming that there exists an integer $j$ such that $1\leq j\leq 2n(\alpha)$ and $E \subset \mathbb{R}^3_{s_j}$ with Lipschitz boundary. We will show that
		\begin{equation}\label{eqn6}
			\dfrac{P_{2,\alpha,j}(B_{s_j})^{3/2}}{|B_{s_j}|_{2,\alpha}} \leq \dfrac{P_{2,\alpha,j}(E)^{3/2}}{|E|_{2,\alpha}}.
		\end{equation}
		
		Without loss of generality, we can assume that $j=1$. We define $\Phi_1: \mathbb{R}_+ \times \left( 0, \dfrac{\pi}{n(\alpha)} \right) \times \mathbb{R} \to \mathbb{R}^3_{s_1}$ by
		$$ \Phi_1 (r, \theta, y) = (r\cos \theta, r\sin \theta, y).$$
		It is easy to see that $\Phi_1$ is a homeomorphism. Next, we define $\Phi_2: \mathbb{R}_+ \times \left( 0, \dfrac{\pi}{n(\alpha)} \right) \times \mathbb{R} \to \mathbb{R}^3_{\widetilde{s}_1}$ by
		$$ \Phi_2 (r, \theta, \eta) = \left(\dfrac{r^{\alpha+1}\cos (\alpha+1)\theta}{\alpha+1}, \dfrac{r^{\alpha+1}\sin(\alpha+1)\theta}{\alpha+1}, \eta\right), $$
		where
		$$\mathbb{R}^3_{\widetilde{s}_1} := \left\{(\xi_1,\xi_2,\eta) = (\rho\cos\varphi, \rho\sin\varphi, y) \in \mathbb{R}^3, \textrm{ } \rho >0, \textrm{ } \varphi \in \Big(0,\frac{(\alpha+1)\pi}{n(\alpha)}\Big) \right\}.$$
		It is also easy to verify that $\Phi_2$ is a homeomorphism. Hence
		$$ \Psi = \Phi_2 \circ \Phi_1^{-1} $$
		is a homeomorphism. 
		
		Next, let $\widetilde{E} = \Psi (E)$ and $\widetilde{B}_{s_1} = \Psi(B_{s_1})$. We deduce that
		$$ \widetilde{B}_{s_1} = \left\{ (\xi_1, \xi_2, \eta) \in \mathbb{R}^3_{\widetilde{s}_1}, \textrm{ } \xi_1^2 + \xi_2^2 + \eta^2 < 1\right\}.$$
		By some straightforward calculations, we obtain
		$$ |E|_{2,\alpha} = \int\limits_{\widetilde{E}} d\xi_1d\xi_2d\eta = |\widetilde{E}|, \quad |B_{s_1}|_{2,\alpha} = \int\limits_{\widetilde{B}_{s_1}} d\xi_1d\xi_2d\eta = |\widetilde{B}_{s_1}|, $$
		$$ P_{2,\alpha,1} (E) = \int\limits_{\partial_{\widetilde{s}_1} \widetilde{E}} d\mathcal{H}^2 =: P_{s_1}(\widetilde{E}), \quad P_{2,\alpha,1} (B_{s_1}) = \int\limits_{\partial_{\widetilde{s}_1} \widetilde{B}_{s_1}} d\mathcal{H}^2 =: P_{s_1}( \widetilde{B}_{s_1}),$$
		where
		$$ \partial_{\widetilde{s}_1} \widetilde{E} = \partial \widetilde{E} \cap \mathbb{R}^3_{\widetilde{s}_1}, \quad \partial_{\widetilde{s}_1} \widetilde{B}_{s_1} = \partial \widetilde{B}_{s_1} \cap \mathbb{R}^3_{\widetilde{s}_1}.$$
		
		Now, by applying \cite[Theorem 1.1]{Lions}, we deduce that
		$$ \dfrac{P_{s_1}(\widetilde{E})^{3/2}}{|\widetilde{E}|} \geq \dfrac{P_{s_1}(\widetilde{B}_{s_1})^{3/2}}{|\widetilde{B}_{s_1}|}, $$
		and thus \eqref{eqn6} follows. 
		
		\noindent \textbf{Step 2:} For the open, bounded set $E \subset \mathbb{R}^3$ with Lipschitz boundary, we set
		$$E = \bigcup_{i=1}^{2n(\alpha)} E_{s_i}, $$
		where $E_{s_i} := E \cap \overline{\mathbb{R}}^3_{s_i}$. It is easy to see that
		\begin{equation}\label{eqn7}
			\bigcup_{i=1}^{2n(\alpha)} \partial_{s_i} E \subset \partial E.
		\end{equation}
		Notice that the sets $E_{s_i}$ are respectively open, bounded subsets in $\overline{\mathbb{R}}^3_{s_i}$ with Lipschitz boundaries. Hence we have the inequality \eqref{eqn6} for each $E_{s_i}$. Since $3/2 > 1$, we have that
		$$ P_{2,\alpha} (E)^{3/2} \geq \sum\limits_{j=1}^{2n(\alpha)} P_{2,\alpha,j}( E_{s_j})^{3/2}.$$
		By combining this and the inequality \eqref{eqn6} for every $E_{s_i}$, our proof is complete. 
	\end{proof}

\begin{thebibliography}{99}
\bibitem[Lions]{Lions} {\sc Lions, P.L., Pacella, F.} \newblock Isoperimetric inequalities for convex cones. \newblock {\em Proc. Amer. Math. Soc.} {\bf 109} (1990), no.~2, 477--485
\end{thebibliography}
\end{document}
